\section{The graph neural network}
\label{part:net}

\subsection{What the model does}

The network learns its own description of a molecule by passing messages along the bonds. Each atom starts from the encoding of \cref{eq:atom-encoder}. In each round, every atom collects a message from each of its neighbours, adds the messages up, and updates itself from the sum. The version used here is the edge-aware graph isomorphism network \parencite{GIN,GINE}, which adds the bond's own encoding into the message, so the kind of bond matters as much as the neighbour on the other end of it:
\begin{equation}
    \begin{split}
        h_i^{(k+1)} = \mathrm{MLP}^{(k)}\Big( &(1+\varepsilon^{(k)})\, h_i^{(k)} \\
        &+ \sum_{j \in \mathcal{N}(i)} \mathrm{ReLU}\big(h_j^{(k)} + e_{ji}\big) \Big).
    \end{split}
    \label{eq:gine}
\end{equation}
That detail is what lets the model tell an aromatic ring from a plain one with the same shape. After three rounds every atom knows about everything within three bonds of it, which covers a chemical group and its surroundings. The atom states are then averaged into one vector for the whole molecule,
\begin{equation}
    h_G = \frac{1}{|V|} \sum_{i \in V} h_i^{(K)},
    \label{eq:pool}
\end{equation}
and a small two-layer head turns that into a single score. Averaging means the answer does not depend on the order the atoms happen to be listed in, but it also means that in a very small molecule each atom has a large say --- a point \ref{part:eval} returns to.

The full model is three such layers, each wrapping a $64 \to 128 \to 64$ perceptron with batch normalisation \parencite{BatchNorm}, ReLU and dropout 0.3, for \num{67588} trainable values in all.

\subsection{How it was trained}

Training used binary cross-entropy with the positive class down-weighted by $n_-/n_+ \approx 0.31$ from the training counts, which is the same balancing the baseline gets, so neither model has an easier objective. The optimiser was Adam \parencite{Adam} at learning rate \num{1e-3}, batches held 64 molecules, and training stopped early when validation ROC-AUC had not improved for 25 epochs, restoring the best checkpoint. The ceiling of 150 epochs was never reached; all six runs stopped between 54 and 100 epochs. Three seeds, 0, 1 and 2, were run for each split, while the split seed stayed at 42, so the effect of the seed and the effect of the split can be told apart.

\subsection{What the curves show}

\Cref{fig:curves} shows all six runs and \cref{tab:seeds} summarises them. Training loss falls smoothly to 0.08--0.12, while validation loss reaches its lowest point around epoch 5--10 and then climbs back to 0.30--0.40. That gap is plain overfitting, which is no surprise with \num{67588} values to fit and 1375 training molecules. Nothing is underfitting: training loss is still falling when training stops.

The interesting part is that validation ROC-AUC does \emph{not} get worse while validation loss climbs. The explanation is that the network grows over-confident, pushing its probabilities towards 0 and 1, and cross-entropy punishes confident mistakes very harshly --- but the \emph{order} it puts the molecules in keeps improving. Since ranking is what ROC-AUC measures, stopping on validation loss would have halted training around epoch 8 and thrown away a better model. This is why the checkpoint rule watches ROC-AUC instead, and why a decision threshold has to be chosen on validation data rather than left at 0.5. The spikes in the loss curve and the sawtooth in the ROC-AUC curve come from the small validation sets of about 295 molecules, where a handful of confident mistakes moves the average sharply; they do not grow over time, so they are noise around a settled model rather than a sign of instability.

Seed and split pull with very different strengths. Across seeds, validation ROC-AUC varies by only $\pm 0.0022$ on the random split and $\pm 0.0025$ on the scaffold split. Changing the split moves the average by 0.0354 --- about \textbf{seven times} as much --- and the two groups do not overlap at all, since the worst random-split seed (0.9130) still beats the best scaffold-split seed (0.8822). The drop is caused by the splitting protocol, not by luck in the initialisation. What the seed does change is the route taken: the best epoch lands anywhere between 28 and 74 across runs that finish at almost the same quality, which is a good argument for checkpointing on a validation score rather than training for a fixed number of epochs.
